\section{The explicit two-orbit placement}
\label{sec:orbit}

The periods \(t,u\) and the half-turn \(R=\half r\) are not merely existence witnesses.  They give a closed formula for every tile in the Conway tiling.

\subsection{The placement family}

Recall
\[
 T^0_{m,n}=T+mt+nu,
 \qquad
 T^1_{m,n}=R(T)+mt+nu.
\]
The superscript records one of two orbits under the translation lattice
\[
 L=\Z t+\Z u.
\]
The complete local adjacency of a direct tile is as follows.

\begin{center}
\begin{tabular}{@{}lll@{}}
\toprule
Boundary path of \(T^0_{m,n}\) & Matching path & Neighboring tile \\
\midrule
\(a\) & \(d\) & \(T^0_{m-1,n}\) \\
\(d\) & \(a\) & \(T^0_{m+1,n}\) \\
\(b\) & \(b\) & \(T^1_{m,n-1}\) \\
\(c\) & \(c\) & \(T^1_{m+1,n-1}\) \\
\(e\) & \(e\) & \(T^1_{m,n}\) \\
\(f\) & \(f\) & \(T^1_{m-1,n}\) \\
\bottomrule
\end{tabular}
\end{center}

For instance, \cref{eq:bpair} says
\[
 T^1_{m,n-1}(b)
 =mt+(n-1)u+\half r(b)
 =-\bigl(mt+nu+b\bigr),
\]
where the minus sign means reversal of orientation, not pointwise negation.  Thus it matches the \(b\)-path of \(T^0_{m,n}\) exactly.  For \(c\), the shift in \cref{eq:cpair} is \(u-t\), which explains the neighbor \((m+1,n-1)\).  The other rows are immediate from \cref{eq:ad,eq:epair,eq:fpair}.

\begin{proposition}[Every open boundary segment is paired]
\label{prop:edgepaired}
In the family \cref{eq:tilefamily}, every oriented boundary path of every tile occurs exactly with its reverse as a boundary path of a tile listed in the adjacency table.  The pairing is involutive.
\end{proposition}
\begin{proof}
The table proves the assertion for direct tiles.  Reversing each row proves it for the corresponding turned tile.  Algebraically, the neighbor maps are
\[
 (0,m,n,a)\leftrightarrow(0,m-1,n,d),
 \quad
 (0,m,n,b)\leftrightarrow(1,m,n-1,b),
\]
and similarly for the other four labels.  Applying any map twice returns the original index.  All equalities are consequences of \cref{lem:orbitpairs} and integer addition.
\end{proof}

\subsection{Why pairwise edge matching is not yet a proof}

\Cref{prop:edgepaired} rules out an exposed open boundary arc, but by itself it does not prove that the plane is covered exactly once.  Around a marked endpoint, paired arcs could conceivably wind so as to leave a gap, or several tiles could overlap.  A picture suggests the correct cyclic order, but a formal proof must derive it.

One classical repair constructs a surface from abstract tile copies and invokes covering-space theory.  That route is valid, but it imports more topology than iPRA needs.  We instead quotient the explicit placement by the period lattice.  Only two tiles remain.  Their sum is an oriented two-cycle on a finite torus subdivision.  The cycle condition forces a globally constant integer multiplicity, and an area identity computes that integer as one.  Thus the same calculation settles all vertices simultaneously.

\subsection{The group generated by the permitted half-turns}

The construction also explains why only the designated half-turns are needed.  The half-turns about doubled centers \(p,q,r,s\) generate the translations
\[
 t=q-p=r-s,
 \qquad
 u=r-p,
\]
by \cref{eq:hh}.  Conversely every turned tile in \cref{eq:tilefamily} is obtained by \(R=\half r\) followed by a lattice translation.  Therefore all placements lie in the group generated by the original boundary half-turns.  This is the wallpaper-group \(p2\) mechanism described geometrically in~\cite{Schattschneider1980}, written here as exact integer-vector algebra.
